\documentclass[11pt,reqno]{amsart}

\usepackage[T1]{fontenc}
\usepackage[utf8]{inputenc}
\usepackage{amsmath,amssymb,amsthm}
\usepackage{array}
\usepackage{booktabs}
\usepackage{xurl}
\usepackage{hyperref}

\emergencystretch=1.8em
\makeatletter
\g@addto@macro\ttfamily{\def\_{\textunderscore\penalty\z@}}
\makeatother

\newtheorem{theorem}{Theorem}[section]
\newtheorem{proposition}[theorem]{Proposition}
\newtheorem{lemma}[theorem]{Lemma}
\newtheorem{corollary}[theorem]{Corollary}
\theoremstyle{remark}
\newtheorem{remark}[theorem]{Remark}

\newcommand{\R}{\mathbb{R}}
\newcommand{\C}{\mathbb{C}}
\newcommand{\Z}{\mathbb{Z}}
\newcommand{\N}{\mathbb{N}}
\newcommand{\B}{\mathbb{B}}
\newcommand{\Hy}{\mathbb{H}}

\title[Atiyah--Sutcliffe Counterexamples and No-Go Rigidity]{Counterexamples and Holomorphic No-Go Rigidity for the Euclidean, Hyperbolic, and Symplectic Atiyah--Sutcliffe Conjectures}
\author{Tomasz Tar\l owski}
\date{November 24, 2026}
\subjclass[2020]{Primary 51M15, 70G45; Secondary 22E46, 81Q70, 68V20}

\begin{document}

\begin{abstract}
We present explicit finite witnesses, continuous families, universal rigidity theorems, and Lean~4 proofs refuting the Atiyah--Sutcliffe conjectures across Euclidean $C_n(\mathbb{R}^3)$, hyperbolic $C_n(\mathbb{H}^3)$, and classical Lie-group loci $\mathrm{Sp}(m)$ and $\mathrm{SO}(2m+1)$.
First, we formalize a Lean~4 disproof of Euclidean Conjecture~3 at the minimal $n = 5$ in $\mathbb{R}^2, \mathbb{R}^3$ ($R_3 \approx 0.999996368278 < 1$; planar failure first announced by Tan--Zhang, Oct.~2026), prove Hyperbolic C3 false at $n = 5$ ($\mathbb{H}^2, \mathbb{H}^3$), and establish an exact $5$-gap positivity identity showing $D_3^{\mathbb{H}} > 1$ on cyclic ideal-slope triangles in $\mathbb{B}^3$.
Second, via certified kernel interval arithmetic and $y$-reflection symmetry, we lower C1 and direct nonzero C2 ($-1 < D_n < 0$) disproofs to planar $n = 35$ ($\mathbb{R}^2, \mathbb{H}^2$), non-coplanar $n = 37$ ($\mathbb{R}^3, \mathbb{H}^3$), and strictly non-coplanar $n = 26$ in both $\mathbb{R}^3$ and $\mathbb{H}^3$ ($D_{26} = 0$, $D_{26}^{\mathbb{H}} = 0$, $-1 < D_{26}^{\mathbb{H}} < 0$, alongside $n = 36, 57, 59, 63$).
Third, bypassing quaternionic Study non-negativity $D_{2m} \ge 0$ via $z$-rotation parity-subblock factorization, we refute C1, C2, and $D \ne 0$ on centrally symmetric planar $\mathrm{Sp}(20)$ ($2m = 40$), 3D $\mathrm{Sp}(18)$ ($2m = 36$), and origin-augmented $B_{18} = \mathrm{SO}(37)$ ($n = 37$, including linear dependence of the $36$ non-origin rows).
Fourth, proving $\mathrm{SU}(2) \times S_m$ multilinear spinor rigidity, we deduce unconditional No-Go theorems showing that arbitrary finite multi-body scalar-weighted sums of equivariant multilinear spinor polynomials lose linear independence at $n = 26$ ($\mathbb{R}^3$), $n = 35$ ($\mathbb{R}^2$), and $2m = 36$ ($\mathrm{Sp}(18)$).
\end{abstract}

\maketitle

\section{Introduction}

Motivated by the Berry--Robbins geometric spin-statistics construction~\cite{BerryRobbins1997} for $n$ identical particles in $\R^3$, Atiyah~\cite{Atiyah2001} and Atiyah--Sutcliffe~\cite{AtiyahSutcliffe2002} associated to every configuration $\mathbf{x} = (x_0, \dots, x_{n-1}) \in C_n(\R^3)$ of $n$ pairwise distinct points an $n$-tuple of complex polynomials $p_0, \dots, p_{n-1} \in \C[t]_{\le n-1}$ of degree at most $n - 1$.
Specifically, for each ordered pair $i \ne j$, let $v_{ij} = (x_j - x_i)/\|x_j - x_i\| \in S^2$, and choose a nonzero $+1$-eigenspinor lift $\psi_{ij} \in \C^2 \setminus \{0\}$ satisfying $(v_{ij} \cdot \boldsymbol{\sigma})\psi_{ij} = \psi_{ij}$, where $\boldsymbol{\sigma} = (\sigma_1, \sigma_2, \sigma_3)$ are the standard Pauli matrices
\[
  \sigma_1 = \begin{pmatrix} 0 & 1 \\ 1 & 0 \end{pmatrix}, \qquad
  \sigma_2 = \begin{pmatrix} 0 & -i \\ i & 0 \end{pmatrix}, \qquad
  \sigma_3 = \begin{pmatrix} 1 & 0 \\ 0 & -1 \end{pmatrix}.
\]
Any two $+1$-eigenspinors of $v_{ij} \cdot \boldsymbol{\sigma}$ differ only by a nonzero complex scalar.
Writing $\omega(u, v) = u_0 v_1 - u_1 v_0$ for the standard $\mathrm{SL}(2, \C)$-invariant symplectic form on $\C^2$, define the Atiyah polynomials and coefficient matrix $M(\mathbf{x}) = (M_{ik}(\mathbf{x}))_{0 \le i,k < n}$ via
\[
  p_i(t) = \prod_{j \ne i} \bigl(\psi_{ij,0} + \psi_{ij,1} t\bigr) = \sum_{k=0}^{n-1} M_{ik}(\mathbf{x}) t^k,
\]
and set the normalized Atiyah--Sutcliffe determinant to be
\begin{equation}\label{eq:def-Dn}
  D_n(\mathbf{x}) := \frac{\det M(\mathbf{x})}{\prod_{0 \le i < j < n} \omega(\psi_{ij}, \psi_{ji})}.
\end{equation}
Because rescaling $\psi_{ij} \mapsto \lambda_{ij}\psi_{ij}$ ($\lambda_{ij} \in \C^\times$) multiplies both numerator and denominator by $\prod_{i \ne j} \lambda_{ij}$, $D_n(\mathbf{x})$ is independent of spinor lift choices and is invariant under point permutations, $\mathrm{SO}(3)$ rotations, translations, and positive dilations~\cite{AtiyahSutcliffe2002}.

Atiyah~\cite{Atiyah2001} and Atiyah--Sutcliffe~\cite{AtiyahSutcliffe2002} posed three central conjectures on $C_n(\R^3)$, together with their hyperbolic $\Hy^3$ analogues obtained by replacing Euclidean unit vectors $v_{ij}$ with forward ideal boundary endpoints $u_{ij} \in S^2 = \partial_\infty \Hy^3$ of oriented hyperbolic geodesic rays from $x_i$ through $x_j$:
\begin{itemize}
  \item \textbf{Conjecture~1 (C1, $n \ge 2$).} For every configuration of $n$ distinct points, $p_0, \dots, p_{n-1}$ are linearly independent; equivalently, $D_n \ne 0$.
  \item \textbf{Conjecture~2 (C2, $n \ge 2$).} For every configuration of $n$ distinct points, $|D_n| \ge 1$.
  \item \textbf{Conjecture~3 (C3, $n \ge 3$).} For every configuration $\mathbf{x}$ of $n$ distinct points,
  \[
    |D_n(\mathbf{x})|^{n-2} \ge \prod_{k=0}^{n-1} |D_{n-1}(\mathbf{x} \setminus \{x_k\})|.
  \]
\end{itemize}
Conjecture~2 directly implies Conjecture~1.
Atiyah and Bielawski~\cite{AtiyahBielawski2002} generalized the Berry--Robbins existence problem from $\mathrm{U}(n)$ to arbitrary compact Lie groups $G$ via Nahm's equations (non-elementary ODEs), while Malkoun~\cite{Malkoun2014,Malkoun2015} constructed explicit Atiyah--Sutcliffe polynomial candidates for all remaining classical groups, notably the symplectic groups $\mathrm{Sp}(m)$ (corresponding to centrally symmetric $2m$-point configurations $\mathbf{X} = -\mathbf{X}$) and the odd orthogonal groups $B_m = \mathrm{SO}(2m+1)$ (centrally symmetric pairs together with a fixed particle at the origin).

In the classical $\mathrm{U}(n)$ Euclidean setting, Eastwood--Norbury~\cite{EastwoodNorbury2001} proved C1 for $n = 4$ (with additional special families proved by \DJ okovi\'c~\cite{Djokovic2002}), Atiyah~\cite{Atiyah2001} proved C3 for $n = 3$, Mazur--Petrenko~\cite{MazurPetrenko2011} established algebraic reformulations and special-case bounds on C1--C3, and Bou~Khuzam--Johnson~\cite{BouKhuzamJohnson2014} established Euclidean C2 and C3 for $n = 4$.
In October~2026, Eickhoff~\cite{Eickhoff2026} proved the asymptotic existence of arbitrarily large finite singular configurations ($D_n = 0$, hence refuting C1 and asymptotically C2) in planar Euclidean space via two-orbit dihedral families ($K = 2$), and established asymptotic hyperbolic singular configurations via small-scale transfer into the Beltrami--Klein ball (\cite[\S7]{Eickhoff2026}), while leaving open concrete small-$n$ bounds, explicit C2 counterexamples with $-1 < D_n < 0$, Conjecture~3, and whether additional orbit rows ($K > 2$) yield new mechanisms or lower concrete thresholds.
Independently in October~2026, Akhtyamov~\cite{Akhtyamov2026} formalized a disproof of Euclidean Conjecture~1 in Lean~4 using a $46$-point helical 3D configuration via the unnormalized projective formulation of~\cite{MazurPetrenko2011} (which expresses C1 without defining the normalized determinant $D_n$ needed for C2 and C3).
Also in October~2026, Tan and Zhang~\cite{TanZhang2026} announced that $\operatorname{Re} D_5 \ge 5/768$ for every five-point configuration (hence C1 holds for $n = 5$) and exhibited a planar integer five-point configuration violating Conjecture~3, certified by rational interval arithmetic, noting that five is the smallest cardinality at which C3 can fail.
Their preprint predates the present work; our five-point witness in Section~\ref{sec:n5} was found without knowledge of it, and we make no priority claim for the failure of Euclidean C3.

Here we formalize the normalized determinant $D_n(\mathbf{x})$ directly in Lean~4 (\texttt{v4.37.0-rc1} + Mathlib, zero \texttt{sorry}, zero \texttt{admit}, zero \texttt{native\_decide}, standard axioms only) and give a unified resolution of both the classical conjectures and Atiyah's broader Lie-theoretic and algebraic program:

\begin{theorem}[Sharp Euclidean C3 disproof at $n = 5$, planar and non-coplanar]\label{thm:main-c3}
Let $\mathbf{x}^{(5)} \in C_5(\R^2) \subset C_5(\R^3)$ be the explicit $5$-point planar integer configuration defined in Section~\ref{sec:n5}.
Then
\[
  |D_5(\mathbf{x}^{(5)})|^{5-2} < \prod_{k=0}^{4} |D_4(\mathbf{x}^{(5)} \setminus \{x_k\})|.
\]
Moreover, lifting $x_0$ off the plane by sufficiently small $\delta > 0$ yields an existentially certified strictly non-coplanar configuration $\mathbf{x}^{(5)}_\delta \in C_5(\R^3)$ violating Conjecture~3.
Since Euclidean C3 holds for $n = 3, 4$, $n = 5$ is the exact minimal particle count for which it fails in $\R^3$.
\end{theorem}

\begin{theorem}[Planar refutations of C1 and C2 down to $n = 35$]\label{thm:main-c12}
Let $\mathbf{x}^{(35)}(s) \in C_{35}(\R^2) \subset C_{35}(\R^3)$, $\mathbf{x}^{(57)}(s) \in C_{57}(\R^2) \subset C_{57}(\R^3)$ ($K = 3$), and $\mathbf{x}^{(63)}(s) \in C_{63}(\R^2) \subset C_{63}(\R^3)$ ($K = 2$) ($s \in [0, 1]$) be the explicit linear homotopies defined in Section~\ref{sec:planar}.
In Lean~4:
\begin{enumerate}
  \item At the explicit integer endpoint $\mathbf{x}^{(35)}(1)$, $D_{35}(\mathbf{x}^{(35)}(1)) \in \R$ satisfies $-1 < D_{35}(\mathbf{x}^{(35)}(1)) < 0$ (numerically $-0.49960620980330087887$, certified at scale $S_{35} = 10^{34}$), directly refuting Conjecture~2; by the Intermediate Value Theorem along $\mathbf{x}^{(35)}$, Lean~4 certifies a planar singular configuration with $D_{35}(\mathbf{x}^{(35)}(s_0)) = 0$ on $(0, 1)$ (numerically $s_0 \approx 0.99999999999999792233$), refuting Conjecture~1 at $n = 35$.
  \item At $n = 57$ ($K=3$) and $n = 63$ ($K=2$), Lean~4 likewise certifies $-1 < D_{57}(\mathbf{x}^{(57)}(1)) < 0$ ($\approx -0.498322440876$), $-1 < D_{63}(\mathbf{x}^{(63)}(1)) < 0$ ($\approx -0.639592897554$), and interior homotopy zeros $D_{57} = 0$ and $D_{63} = 0$ on $(0, 1)$.
\end{enumerate}
(Here the exported existential declarations \texttt{exists\_planar\_35\_atiyahDet\_zero}, \texttt{exists\_planar\_57\_atiyahDet\_zero}, and \texttt{exists\_planar\_63\_atiyahDet\_zero} state the existence of planar singular configurations in $C_n(\R^3)$, with interiority $s_0 \in (0, 1)$ along each path following from their IVT proofs via $D_n(\mathbf{x}(0)) = 1$ and $D_n(\mathbf{x}(1)) < 0$.)
\end{theorem}

\begin{theorem}[Strictly non-coplanar 3D refutations of C1 and C2 at $n = 26$, $36$, $37$, and $59$]\label{thm:main-3d}
In strictly non-coplanar Euclidean three-space, Lean~4 certifies four families refuting Conjectures~1 and~2:
\begin{enumerate}
  \item At $n = 26$, the explicit mirror-symmetric integer family $\mathbf{x}^{(26)}(t) \in C_{26}(\R^3)$ (\texttt{x26Clamp}, Section~\ref{sec:3d26}) satisfies non-coplanarity (\texttt{Non\allowbreak coplanar3D}) for every $t \in \R$ and $\mathrm{Im}\, D_{26}(\mathbf{x}^{(26)}(t)) = 0$ identically, with a Lean-certified sign reversal $D_{26}(\mathbf{x}^{(26)}(8)) > 0$ and $D_{26}(\mathbf{x}^{(26)}(10)) < 0$ via complex interval arithmetic at scale $S_{26} = 10^{30}$. Continuity yields an interior zero $t_0 \in (8, 10)$ with $D_{26}(\mathbf{x}^{(26)}(t_0)) = 0$ (numerically $t_0 \approx 8.289031$), refuting both C1 and C2 at $n = 26$.
  \item At $n = 36$, the coaxial two-ring Pythagorean family $\mathbf{x}^{(36)}(t) \in C_{36}(\R^3)$ (\texttt{x36Clamp}, $(33, 56, 65)$ rings, Section~\ref{sec:3d36}) is strictly non-coplanar for every $t \in \R$ and satisfies $\mathrm{Im}\, D_{36}(\mathbf{x}^{(36)}(t)) = 0$ identically, with a Lean-certified sign reversal $D_{36}(\mathbf{x}^{(36)}(9)) > 0$ and $D_{36}(\mathbf{x}^{(36)}(10)) < 0$ at scale $S_{36} = 10^{18}$, giving an interior zero $D_{36}(\mathbf{x}^{(36)}(t_0)) = 0$ on $(9, 10)$ (numerically $t_0 \approx 9.392050$).
  \item At $n = 37$ (lowering our mirror-pair $n = 59$ direct nonzero C2 bound), appending $(0, \pm 1, H)$ to $\mathbf{x}^{(35)}(s)$ (or to $\mathbf{x}^{(57)}(s)$ for $n = 59$) preserves $\mathrm{Im}\, D_{n+2} = 0$, enforces strict non-coplanarity ($\neg \mathrm{IsCoplanar}$) for $s \ne 0$, and converges to $D_n(\mathbf{x}^{(n)}(s))$ as $H \to +\infty$. Consequently, for sufficiently large $H > 0$, Lean~4 certifies strictly non-coplanar configurations in $C_{37}(\R^3)$ (and $C_{59}(\R^3)$) satisfying both $-1 < D < 0$ (direct nonzero C2 violation) and $D = 0$ (C1 violation).
\end{enumerate}
\end{theorem}

\begin{theorem}[Hyperbolic $\Hy^3$ disproofs ($n = 5, 26, 35, 37$) and $n = 3$ $5$-gap positivity]\label{thm:main-hyp}
For every Euclidean configuration $\mathbf{x} \in C_n(\R^3)$ embedded into the Beltrami--Klein ball $\B^3 \cong \Hy^3$, the tangent scaling limit $\lim_{\varepsilon \to 0^+} D_n^{\Hy}(\varepsilon \mathbf{x}) = D_n(\mathbf{x})$ holds while preserving geodesic vs.\ affine coplanarity.
Lean~4 certifies:
\begin{enumerate}
  \item Hyperbolic Conjecture~3 fails at $n = 5$ in both planar $\Hy^2$ and strictly non-coplanar $\Hy^3$.
  \item Hyperbolic Conjectures~1 and~2 fail in planar $\Hy^2 \subset \Hy^3$ at $n = 35$ (and $n = 57$) and in strictly non-coplanar $\Hy^3$ at $n = 37$ (and $n = 59$) via both direct nonzero witnesses ($-1 < D_n^{\Hy} < 0$) and vanishing configurations ($D_n^{\Hy} = 0$).
  \item Generalizing reality symmetry to arbitrary $y$-reflecting permutations transfers \texttt{x26Clamp} into $\B^3$, lowering the Lean-certified strictly non-coplanar $\Hy^3$ C1/C2 disproof to $n = 26$, yielding both $D_{26}^{\Hy} = 0$ and direct nonzero $-1 < D_{26}^{\Hy} < 0$ in $C_{26}(\B^3)$.
  \item At $n = 3$, any equatorial hyperbolic triangle in $y = 0 \subset \B^3$ avoiding $(0,0,-1)$ whose ideal endpoint stereographic slopes satisfy the cyclic $5$-gap chain $t_{20} < t_{10} < t_{12} < t_{02} < t_{01} < t_{21}$ satisfies $\mathrm{Im}\, D_3^{\Hy} = 0$ and $D_3^{\Hy} > 1$ identically, witnessed unconditionally by $\{(-3/5,0,0), (3/5,0,0), (0,0,4/5)\} \subset \B^3$ ($D_3^{\Hy} \approx 1.008290730761$).
\end{enumerate}
Whether Hyperbolic Conjecture~3 can fail at $n = 4$ remains open.
\end{theorem}

\begin{theorem}[Symplectic $\mathrm{Sp}(m)$ and odd-orthogonal $B_m = \mathrm{SO}(2m+1)$ disproofs]\label{thm:main-symp}
Even when restricted to centrally symmetric configurations $\mathbf{X} = -\mathbf{X}$ (Malkoun's symplectic locus~\cite{Malkoun2014}, where quaternionic Study determinants satisfy $D_{2m} \ge 0$ identically) or origin-augmented centrally symmetric configurations (Malkoun's $B_m = \mathrm{SO}(2m+1)$ locus~\cite{Malkoun2015}), Lean~4 certifies via $z$-rotation $\rho_z(x,y,z) = (-x,-y,z)$ parity-subblock factorization and complex interval Gaussian elimination ($S_{\mathrm{symp}} = 10^{28}$):
\begin{enumerate}
  \item Planar symplectic $\mathrm{Sp}(20)$ ($2m = 40$): there exists a planar centrally symmetric $40$-point configuration with $D_{40} = 0$, refuting \texttt{Symplectic\allowbreak AtiyahDetNonzeroProp 20}, \texttt{Symplectic\allowbreak AtiyahConjecture1Prop 20}, and \texttt{Symplectic\allowbreak AtiyahSutcliffe\allowbreak Conjecture2Prop 20}.
  \item Strictly non-coplanar 3D symplectic $\mathrm{Sp}(18)$ ($2m = 36$): there exists a strictly non-coplanar centrally symmetric $36$-point configuration with $D_{36} = 0$, refuting all three symplectic conjectures at $m = 18$.
  \item Strictly non-coplanar 3D orthogonal $B_{18} = \mathrm{SO}(37)$ ($n = 37$): there exists an origin-augmented centrally symmetric non-coplanar $37$-point configuration with $D_{37} = 0$ whose $36$ non-origin Atiyah polynomials are already linearly dependent, refuting all three $B_{18}$ conjectures.
\end{enumerate}
\end{theorem}

\begin{theorem}[Universal $\mathrm{SU}(2) \times S_m$ holomorphic spinor rigidity and multi-body No-Go]\label{thm:main-nogo}
Every degree-$\le m$ polynomial map $\Phi : (\C^2)^m \to \C[t]$ that is coordinatewise $\C$-linear, $S_m$-symmetric, and $\mathrm{SU}(2)$-equivariant under the standard homogeneous binary action satisfies
\[
  \Phi(v_0, \dots, v_{m-1}) = c \prod_{k=0}^{m-1} (v_{k,0} + v_{k,1} t)
\]
for some universal scalar $c \in \C$ (Lean~4 declaration \texttt{su2\_symmetric\_multilinear\_eq\_atiyahFactor}).
Consequently, for every finite index set $\iota$, every family $(\Phi_\alpha)_{\alpha \in \iota}$ of multilinear $\mathrm{SU}(2) \times S_{n-1}$-equivariant spinor polynomial maps, and completely arbitrary scalar weights
\[
  W_\alpha : C_n(\R^3) \times \{0, \dots, n - 1\} \to \C
\]
(including arbitrary $\mathrm{E}(3)$-invariant multi-body radial weights), the weighted candidate $P_i(\mathbf{x}; t) = \sum_{\alpha \in \iota} W_\alpha(\mathbf{x}, i) \Phi_\alpha((\psi_{ij})_{j \ne i})$ collapses row-wise to $\mu_i(\mathbf{x}) p_i^{\mathrm{Atiyah}}(t)$ and fails to be linearly independent on:
\begin{enumerate}
  \item a strictly non-coplanar 3D configuration at $n = 26$ (\texttt{universal\_holomorphic\_spinor\_nogo\_noncoplanar\_3d\_26});
  \item a planar configuration at $n = 35$ (\texttt{universal\_holomorphic\_spinor\_nogo\_planar\_35});
  \item a centrally symmetric strictly non-coplanar 3D $\mathrm{Sp}(18)$ configuration at $n = 36$ (\texttt{universal\_holomorphic\_spinor\_nogo\_symplectic\_36}).
\end{enumerate}
\end{theorem}

\begin{remark}[Scope relative to Malkoun's graph generalization]
Malkoun~\cite{Malkoun2025} introduced a graph-theoretic extension $A_G(\mathbf{x})$ for arbitrary simple graphs $G$ on $n$ vertices (with $G = K_n$ recovering $D_n$ up to normalization).
This paper and the seven Track~B Lean~4 modules documented in Section~9 treat the complete-graph conjectures on $C_n(\R^3)$ and $C_n(\Hy^3)$, their classical Lie-group loci ($\mathrm{Sp}(m)$, $\mathrm{SO}(2m+1)$), and holomorphic spinor No-Go rigidity.
\end{remark}

\section{Ordered Planar Formula and Chart-Independent Limit Theorem}

Restrict attention first to configurations in the $(x, z)$-plane $y = 0$ with strictly increasing vertical coordinates $z_0 < z_1 < \dots < z_{n-1}$, writing $\mathbf{x}_i = (x_i, 0, z_i)$.
For every $0 \le i < j < n$, let $\Delta x_{ij} = x_j - x_i$, $\Delta z_{ij} = z_j - z_i > 0$, and $r_{ij} = \sqrt{\Delta x_{ij}^2 + \Delta z_{ij}^2}$, and define the stereographic planar slope
\begin{equation}\label{eq:planar-slope}
  c_{ij} := \frac{\Delta x_{ij}}{\Delta z_{ij} + r_{ij}} \in (-1, 1).
\end{equation}
Notice that $v_{ij} = \bigl(\frac{2c_{ij}}{1 + c_{ij}^2}, 0, \frac{1 - c_{ij}^2}{1 + c_{ij}^2}\bigr)$.
Direct substitution into $v_{ij} \cdot \boldsymbol{\sigma}$ shows that $(1, c_{ij})^\top$ is a nonzero $+1$-eigenspinor of $v_{ij} \cdot \boldsymbol{\sigma}$ and $(-c_{ij}, 1)^\top$ is a nonzero $+1$-eigenspinor of $(-v_{ij}) \cdot \boldsymbol{\sigma}$ for every $c_{ij} \in (-1, 1)$ (including the collinear vertical case $c_{ij} = 0$, where $\psi^{\mathrm{can}}(0, 0, -1)$ has a south-pole division singularity).
Moreover, $\omega((1, c_{ij})^\top, (-c_{ij}, 1)^\top) = 1 + c_{ij}^2 > 0$.

\begin{proposition}[Smooth ordered-height planar determinant]\label{prop:planar}
Let $\mathbf{x}_i = (x_i, 0, z_i) \in \R^3$ satisfy $z_0 < z_1 < \dots < z_{n-1}$, and define $M^{\mathrm{pl}}(\mathbf{x}) \in \R^{n \times n}$ by letting row $i$ hold the coefficients $1, t, \dots, t^{n-1}$ of
\[
  p_i^{\mathrm{pl}}(t) = \prod_{j < i} (-c_{ji} + t) \prod_{j > i} (1 + c_{ij} t).
\]
Then $D_n(\mathbf{x})$ is real and equals
\[
  D_n(\mathbf{x}) = \frac{\det M^{\mathrm{pl}}(\mathbf{x})}{\prod_{0 \le i < j < n} (1 + c_{ij}^2)}.
\]
In particular, along any path $s \mapsto \mathbf{x}(s) = ((s X_i, 0, Z_i))_{i=0}^{n-1}$ with $Z_0 < \dots < Z_{n-1}$, every point remains distinct for all $s \in \R$, $s \mapsto D_n(\mathbf{x}(s))$ is real-valued and continuous, and $D_n(\mathbf{x}(0)) = 1$ since $c_{ij}(0) = 0$ makes $M^{\mathrm{pl}}(\mathbf{x}(0)) = I_n$.
(Lean~4: \texttt{atiyahDet\_eq\_planarMat\_div}, \texttt{continuous\_planarPath\_re}, \texttt{planarPath\_zero\_eq\_one}.)
\end{proposition}

\begin{proof}
Since $\Delta z_{ij} > 0$, $\Delta z_{ij} + r_{ij} > 0$ avoids zero denominators across all $0 \le i < j < n$.
By spinor-lift independence (Lemma~\ref{lem:chart-limit}), choosing $\psi_{ij} = (1, c_{ij})^\top$ for $i < j$ and $\psi_{ij} = (-c_{ji}, 1)^\top$ for $i > j$ produces rows equal to $p_i^{\mathrm{pl}}(t)$ and denominator $\prod_{i<j} (1 + c_{ij}^2) > 0$, establishing reality and continuity across all $s \in \R$.
\end{proof}

To prove genuine 3D and hyperbolic convergence where unit directions $v_{ij}$ can approach the south pole $(0, 0, -1)$ (discontinuous for the normalized northern formula), we formalize local polynomial chart switching:

\begin{lemma}[Chart-independent spinor determinant and unit-direction limit]\label{lem:chart-limit}
Let $D_n^{\mathrm{spin}}(\psi)$ denote the determinant ratio~\eqref{eq:def-Dn} evaluated on any family of nonzero $+1$-eigenspinors $\psi_{ij}$ of $u_{ij} \cdot \boldsymbol{\sigma}$, where $u_{ij} \in S^2$ satisfies $u_{ij} \ne u_{ji}$ for all $i \ne j$.
Then $D_n^{\mathrm{spin}}(\psi)$ is independent of the choice of $+1$-eigenspinors (\texttt{atiyahDetOfSpinors\_eq\_of\_unitSpinors}), so $D_n^{\mathrm{spin}}(u)$ is well defined purely as a function of the direction field $u = (u_{ij})_{i \ne j}$.
For a fixed reference unit vector $u_0 \in S^2$ and arbitrary $w \in S^2$, define the unnormalized polynomial spinor lift
\[
  \psi_{u_0}^{\mathrm{chart}}(w) := \begin{cases}
    (1 + w_2,\; w_0 + i w_1)^\top & \text{if } u_{0,2} \ge 0,\\
    (w_0 - i w_1,\; 1 - w_2)^\top & \text{if } u_{0,2} < 0.
  \end{cases}
\]
Direct multiplication verifies $(w \cdot \boldsymbol{\sigma})\psi_{u_0}^{\mathrm{chart}}(w) = \psi_{u_0}^{\mathrm{chart}}(w)$ (\texttt{pauliDot\_chartSpinor}), and $\psi_{u_0}^{\mathrm{chart}}(w) \ne 0$ in a neighborhood of $u_0$ (\texttt{chartSpinor\_ne\_zero}).
Consequently, whenever unit-sphere direction fields $u_{ij}(\alpha) \to u_{ij}^* \in S^2$ along a filter on $\alpha$, with $u_{ij}^* \ne u_{ji}^*$ at the limit and eventually $u_{ij}(\alpha) \ne u_{ji}(\alpha)$ for all $i \ne j$, we have $D_n^{\mathrm{spin}}(u(\alpha)) \to D_n^{\mathrm{spin}}(u^*)$ (\texttt{tendsto\_atiyahDetOfSpinors\_unit}, \texttt{tendsto\_atiyahDet\_of\_tendsto}).
\end{lemma}

\section{Sharp Counterexample to Euclidean Conjecture~3 at $n = 5$}\label{sec:n5}

\subsection{Explicit integer coordinates and high-precision values}
Define the $5$-point planar witness $\mathbf{x}^{(5)} = (p_0, p_1, p_2, p_3, p_4) \in (\Z \times \{0\} \times \Z)^5$ by
\begin{align*}
  p_0 &= (1198773, 0, -320821), &
  p_1 &= (212773, 0, -58745), &
  p_2 &= (-55, 0, 3), \\
  p_3 &= (-56, 0, 4), &
  p_4 &= (-697578, 0, 117714).
\end{align*}
Its $z$-coordinates are strictly increasing and its minimal inter-point distance is $\|p_3 - p_2\| = \sqrt{2} > 0$.
Table~\ref{tab:n5-values} displays $80$-digit \texttt{mpmath} evaluations (rounded to $20$ decimal places) showing $R_3 = |D_5|^3 / \prod_k |D_4(\mathbf{x}^{(5)} \setminus \{p_k\})| \approx 0.999996368278 < 1$.

\begin{table}[h]
\centering
\small
\caption{Planar $5$-point witness $\mathbf{x}^{(5)}$ and its $4$-point sub-determinants ($20$ digits shown from $80$-digit \texttt{mpmath} evaluation).}
\label{tab:n5-values}
\begin{tabular}{ll}
\toprule
Quantity & Numerical value \\
\midrule
$D_5(\mathbf{x}^{(5)})$ & $1.00098008192048394175$ \\
$D_4(\mathbf{x}^{(5)} \setminus \{p_0\})$ & $1.00047373646488409028$ \\
$D_4(\mathbf{x}^{(5)} \setminus \{p_1\})$ & $1.00051482265799591186$ \\
$D_4(\mathbf{x}^{(5)} \setminus \{p_2\})$ & $1.00097400674265848250$ \\
$D_4(\mathbf{x}^{(5)} \setminus \{p_3\})$ & $1.00097394960174767561$ \\
$D_4(\mathbf{x}^{(5)} \setminus \{p_4\})$ & $1.00000711481910285281$ \\
\midrule
LHS $= |D_5(\mathbf{x}^{(5)})|^3$ & $1.00294312838459245275$ \\
RHS $= \prod_{k=0}^4 |D_4(\mathbf{x}^{(5)} \setminus \{p_k\})|$ & $1.00294677080822041528$ \\
Ratio $R_3 = \mathrm{LHS}/\mathrm{RHS}$ & $0.99999636827822375130 < 1$ \\
\bottomrule
\end{tabular}
\end{table}

\subsection{Denominator cancellation and Lean~4 interval certification}
Writing $\mathrm{den}_m(\mathbf{y}) := \prod_{0 \le i < j < m}(1 + c_{ij}(\mathbf{y})^2) > 0$, each pair $\{i, j\}$ ($0 \le i < j < 5$) belongs to exactly $3$ of the $5$ four-point deletions, giving the exact identity $\mathrm{den}_5(\mathbf{x})^3 = \prod_{k=0}^{4} \mathrm{den}_4(\mathbf{x} \setminus \{x_k\})$.
Thus C3 at $n = 5$ reduces cleanly to $|\det M^{\mathrm{pl}}(\mathbf{x}^{(5)})|^3 \ge \prod_{k=0}^{4} |\det M^{\mathrm{pl}}(\mathbf{x}^{(5)} \setminus \{p_k\})|$.
Using outward-rounded fixed-point intervals \texttt{FpIv} at scale $S_5 = 10^{18}$, Lean's kernel certifies an integer separation margin exceeding $3.09 \times 10^{90}$ (\texttt{c3\_int\_gap}, \texttt{exists\_five\_planar\_points\_violating\_conjecture\_3}, \texttt{disproof\_atiyah\_sutcliffe\_conjecture\_3}).

\subsection{Strictly non-coplanar 3D realization at $n = 5$}\label{sec:n5-3d}
Because $p_1, p_2, p_3$ are non-collinear in $y = 0$, lifting $p_0$ off the plane to $(1198773, \delta, -320821)$ for $\delta \ne 0$ (\texttt{x5Perturb}) produces a strictly non-coplanar configuration (\texttt{not\_isCoplanar\_x5Perturb}).
Applying Lemma~\ref{lem:chart-limit} as $\delta \to 0^+$ to $D_5(\mathbf{x}^{(5)}_\delta)$ and all five sub-determinants proves the strict non-coplanar theorem \texttt{exists\_five\_noncoplanar\_points\_violating\_conjecture\_3} of Theorem~\ref{thm:main-c3}.

\section{Planar Records: $n = 63$ ($K=2$), $n = 57$ ($K=3$), and Staggered $n = 35$}\label{sec:planar}

\subsection{Exact $2 \times 2$ axial reduction at $n = 63$ and direct $|D_{63}| < 1$ certificate}\label{sec:n63}
Take $a = 31$, $b = 30$, $n = 63$, and define $\mathbf{x}^{(63)}(s) = ((s X_i, 0, Z_i))_{i=0}^{62}$ by $(0, i)$ for $0 \le i < 31$, $(3097, 38539)$ at $i = 31$, $(80093, 40720)$ at $i = 32$, and $(0, 40830 + i - 33)$ for $33 \le i < 63$.
Because only rows $31$ and $32$ are non-axial, eliminating lower axial pivots left-to-right and upper axial pivots right-to-left reduces $\det M^{\mathrm{pl}}(\mathbf{x}^{(63)}(1))$ identically to a $2 \times 2$ determinant (\texttt{planarMat\_det\_eq\_elimFront}).
At $S_{63} = 10^{35}$, \texttt{FpIv} encloses $S_{63} \det M^{\mathrm{pl}}(\mathbf{x}^{(63)}(1))$ inside $[\mathrm{lo}_{63}, \mathrm{hi}_{63}]$ with $\mathrm{hi}_{63} < 0$, establishing $D_{63}(\mathbf{x}^{(63)}(1)) < 0$ and a homotopy zero $D_{63}(\mathbf{x}^{(63)}(s_0)) = 0$ (numerically $s_0 \approx 0.98861368480927$, $D_{63}(1) \approx -0.639592897554$).

In \texttt{AtiyahSutcliffeExtensions.lean}, we add recursive interval denominator folding \texttt{denFoldIv} over off-axis pairs (\texttt{denFoldIv\_lo\_le\_planarDen}): since $1 + c_{ij}^2 = 1$ identically whenever $x_i = x_j = 0$, only pairs touching $\{31, 32\}$ contribute (\texttt{activePairs\_valid}).
Evaluating at $S_{63} = 10^{35}$ yields exact integer bounds \texttt{den63\_cert} spanning
{\scriptsize
\begin{align*}
  \mathrm{den}_{63}.\mathrm{lo} &= 3572463025754843646294686642028955894118226387223898, \\
  \mathrm{den}_{63}.\mathrm{hi} &= 3572463025754843646294686642028962394159364947130643,
\end{align*}
}%
certifying $|D_{63}(\mathbf{x}^{(63)}(1))| < 1$ directly inside Lean~4 (\texttt{x63Path\_one\_norm\_lt\_one}).

\subsection{Certified Gaussian elimination and the $n = 57$ ($K=3, m=1$) realization}\label{sec:n57}
To certify multi-orbit matrices generally, we formalize outward-rounded integer interval division \texttt{FpIv.divFp} and pivot elimination \texttt{gaussDetIv}: whenever the leading pivot lies in a zero-separated interval (\texttt{SepZero}), $\det M = M_{00} \cdot \det M'$ where $M'_{ij} = M_{i+1,j+1} - (M_{i+1,0}/M_{00}) M_{0,j+1}$ (\texttt{det\_eq\_head\_mul\_det\_elimSubMatrix}, \texttt{gaussDetIv\_mem}).

Define the $57$-point three-orbit planar integer family $\mathbf{x}^{(57)}(s)$ with coordinates $((s X_i, 0, Z_i))_{i=0}^{56}$ by taking $a = 28$ lower axial points $(0, 0, i)$ ($0 \le i < 28$), $K = 3$ strictly positive-$x$ off-axis points
\[
  (14650, 0, 177472), \qquad (418904, 0, 185651), \qquad (457123, 0, 188492),
\]
and $b = 26$ upper axial points $(0, 0, 189224 + i - 31)$ ($31 \le i < 57$).
Running \texttt{polyRowFastIv}, \texttt{gaussDetIv}, and \texttt{denFoldIv} inside Lean's kernel at scale $S_{57} = 10^{40}$ yields exact integer bounds \texttt{det57\_cert} on $S_{57}\det M^{\mathrm{pl}}$ and \texttt{den57\_cert} on $S_{57}\mathrm{den}_{57}$ satisfying $-\mathrm{den}_{57}.\mathrm{lo} < \mathrm{det}_{57}.\mathrm{lo} \le \mathrm{det}_{57}.\mathrm{hi} < 0$, directly certifying $-1 < D_{57}(\mathbf{x}^{(57)}(1)) < 0$ (\texttt{x57Path\_one\_re\_neg\_and\_norm\_lt\_one}, \texttt{exists\_planar\_57\_violating\_conjecture\_2}) and locating a homotopy zero $D_{57}(\mathbf{x}^{(57)}(s_0)) = 0$ (\texttt{exists\_planar\_57\_atiyahDet\_zero}).
Independent $80$-digit \texttt{mpmath} evaluation confirms $D_{57}(\mathbf{x}^{(57)}(1)) \approx -0.49832244087588239334$ and $s_0 \approx 0.99985683874856118709$.

\subsection{Staged Gaussian elimination and the $n = 35$ planar record}\label{sec:n35}
Projecting the staggered-$z$ mirror-pair construction onto the $(x, z)$-plane lowers the planar C1 and C2 threshold dramatically from $n = 57$ to $n = 35$ (\texttt{AtiyahSutcliffePlanar35.lean}).
Specifically, in \texttt{data/atiyah\_part2\_witnesses.json}, the compact $9$-digit integer configuration \texttt{planar\_35\_compact} places $6$ lower axial points $(0, 0, k)$ ($0 \le k \le 5$), $13$ staggered opposite-$x$ pairs $(-X_r, 0, Z_r), (X_r, 0, Z_r + 1)$ ($0 \le r < 13$), and $3$ upper axial points $(0, 0, 508328426 + j)$ ($j = 0, 1, 2$), achieving $D_{35} \approx -7.225271903803 \times 10^{12} < 0$ in $65$-digit \texttt{mpmath}.

To certify $-1 < D_{35} < 0$ and $D_{35} = 0$ simultaneously along one integer homotopy $\mathbf{x}^{(35)}(s) = ((s\,\texttt{pts35X}_i, 0, \texttt{pts35Z}_i))_{i=0}^{34}$ (\texttt{x35Path}), we rescale $Z$ by $10^9$ and contract horizontal coordinates near the crossing point to obtain strictly increasing $z$-coordinates: $6$ lower axial points $(0,0,10^9 k)$ ($0 \le k \le 5$), $13$ staggered pairs $(\mp X_r, 0, Z_{2r}), (\pm X_r, 0, Z_{2r}+10^9)$ at indices $6..31$ with base heights $4769344 \cdot 10^9 \le Z_{2r} \le 371342505 \cdot 10^9$ and horizontal magnitudes $X_r \in [6303169394383633,\, 484559924823925628]$, and $3$ upper axial points $(0,0, (508328426 + j) \cdot 10^9)$ ($j = 0, 1, 2$).
Because direct monolithic kernel evaluation of $35 \times 35$ Gaussian elimination exceeds practical kernel reductions, \texttt{AtiyahSutcliffePlanar35.lean} decomposes row-polynomial construction (\texttt{c35Lit}), elimination pivots (\texttt{Q35\_0} $\dots$ \texttt{Q35\_34}), and denominator folding (\texttt{den35\_cert}) into explicit step traces verified by \texttt{checkTraceLastFp} via kernel-only \texttt{decide +kernel} at scale $S_{35} = 10^{34}$, yielding the exact integer bounds (wrapped across two lines, in units of $10^{-34}$):
{\scriptsize
\begin{align*}
  \texttt{Q35\_34}.\mathrm{lo} &= -58342790689760695113653664093684980097012271 \\
    &\qquad 93778911538203328476262170059640870873735565, \\
  \texttt{Q35\_34}.\mathrm{hi} &= -58342529956444782503259353639093133547143278 \\
    &\qquad 53320120434822824326001383186931254848321828, \\
  \texttt{den35\_cert}.\mathrm{lo} &= \phantom{-}11677729215189810865549331874817352283376480 \\
    &\qquad 683078878412907104944070361450376458840901615, \\
  \texttt{den35\_cert}.\mathrm{hi} &= \phantom{-}11677729215189810865549331874818239576022953 \\
    &\qquad 204200192972936468239435740561964559577426109.
\end{align*}
}%
Decidable inequality checks (\texttt{cert35\_bounds}) yield $\texttt{Q35\_34}.\mathrm{hi} < 0$ and $-\texttt{den35\_cert}.\mathrm{lo} < \texttt{Q35\_34}.\mathrm{lo}$.
This certifies $-1 < D_{35}(\mathbf{x}^{(35)}(1)) < 0$ (\texttt{x35Path\_one\_re\_neg\_and\_norm\_lt\_one}, \texttt{exists\_planar\_35\_violating\_conjecture\_2}) and, by continuity from $D_{35}(\mathbf{x}^{(35)}(0)) = 1$, yields a singular point $s_0 \in (0, 1)$ where $D_{35}(\mathbf{x}^{(35)}(s_0)) = 0$ (\texttt{exists\_planar\_35\_atiyahDet\_zero}).
Independent $65$-digit \texttt{mpmath} computation gives $D_{35}(\mathbf{x}^{(35)}(1)) \approx -0.49960620980330087887$ and root $s_0 \approx 0.99999999999999792233$.

\begin{remark}[Multi-orbit planar comparison ($K=3$ vs.\ $K=2$)]
In~\cite{Eickhoff2026}, Eickhoff analyzed two coaxial $m$-gons ($K = 2$) and noted in the conclusion that additional orbit rings beyond $K = 2$ might generate further singularity families.
Even restricting to $m = 1$ planar configurations with all off-axis points on $x > 0$, our $K = 3$ witnesses achieve negative $D_{57}$ ($D_{57} \approx -1095.557$ for $(a,b) = (27,27)$ in \texttt{data/atiyah\_sutcliffe\_scan.json}, and $-0.498322$ in $\mathbf{x}^{(57)}(1)$), improving by $6$ points on our $K = 2, m = 1$ realization ($n = 63$), while the staggered two-sided $n = 35$ path lowers the certified planar threshold by another $22$ points.
With $K = 2, m = 2$, the $58$-point planar integer family $\mathbf{x}^{(58)}(s)$ numerically satisfies $D_{58}(\mathbf{x}^{(58)}(1)) \approx -0.505421925599$ with root $s_0 \approx 0.999969508354$.
\end{remark}

\section{Strictly Non-Coplanar 3D Disproofs ($n = 26, 36, 37, 59$) and Sweeps}\label{sec:3d}

\subsection{Lean-certified Pythagorean two-ring 3D disproof at $n = 36$}\label{sec:3d36}
To lower the Lean-certified strictly non-coplanar 3D threshold directly without mirror lifting (improving by $10$ points on the $46$-point helical C1 disproof of~\cite{Akhtyamov2026} while certifying normalized $D_{36} = 0$ for C2), \texttt{AtiyahSutcliffe3D36.lean} constructs the explicit $36$-point two-ring Pythagorean family $\mathbf{x}^{(36)}(t) \in C_{36}(\R^3)$ (\texttt{x36Clamp}), rescaling base integer points horizontally by $\max(t, 1)$ while holding heights $z_i$ fixed:
\begin{itemize}
  \item $a = 18$ lower axial points $(0, 0, 10k)$ ($0 \le k \le 17$);
  \item an inner $6$-point Pythagorean $(33, 56, 65)$ ring at $z = 104040$ built from $(520, 264, 448) = 8(65, 33, 56)$, with horizontal coordinates $(\pm 520, 0)$ and $(\pm 264, \pm 448)$ lying on a common circle;
  \item an outer $6$-point Pythagorean ring at $z = 107780$ scaled by ring ratio $R = 23$, with horizontal coordinates $(\pm 11960, 0)$ and $(\pm 6072, \pm 10304)$; and
  \item $b = 6$ upper axial points $(0, 0, 107790 + 10k)$ ($0 \le k \le 5$).
\end{itemize}
The Lean formalization verifies four structural components:
\begin{enumerate}
  \item \textbf{Non-coplanarity, reality, and continuity:} $(0,0,0), (0,0,10)$ and two inner-ring points span $\R^3$ for every $t \in \R$ (\texttt{x36Clamp\_noncoplanar}), while the $y$-reflection involution \texttt{reflectPerm36} forces $\mathrm{Im}\, D_{36}(\mathbf{x}^{(36)}(t)) = 0$ identically (\texttt{atiyahDet\_x36Clamp\_im\_zero}).
  \item \textbf{Antipodal parity block reduction $36 \to 18 \times 18$:} Pairing axial and ring partners under \texttt{antipodalPerm36} negates stereographic slopes across the pair and splits the coefficient matrix via \texttt{parityEquiv} into even- and odd-degree blocks (\texttt{det\_stereoMat\_eq\_even\_odd}): $\det M^{\mathrm{st}}(\mathbf{x}^{(36)}(t)) = 64 \det M_{\mathrm{even}}(t) \det M_{\mathrm{odd}}(t)$.
  \item \textbf{Exact axial elimination $18 \to 6$:} In both \texttt{evenMat18} and \texttt{oddMat18}, axial rows eliminate identically to a $6 \times 6$ Schur complement \texttt{reduce18To6C} (\texttt{det\_18\_eq\_det\_6}).
  \item \textbf{Complex interval sign change and IVT zero:} Outward-rounded complex fixed-point intervals \texttt{CFpIv} at scale $S_{36} = 10^{18}$ (with $9$-stage polynomial checkpoints) certify $D_{36}(9) > 0 > D_{36}(10)$ (80-digit values $+4.4639441897 \times 10^{11}$ and $-1.3361590585 \times 10^{11}$, root $t_0 \approx 9.392050$), proving \texttt{exists\_noncoplanar\_3d\_36\_atiyahDet\_zero}.
\end{enumerate}

\subsection{Lean-certified mirror-symmetric 3D disproof at $n = 26$}\label{sec:3d26}
Releasing the cyclic symmetry while keeping only the reflection $y \mapsto -y$ (which still forces $D_n \in \R$) and alternating local minimization of $\mathrm{Re}\, D_n$ with point deletions and exchanges (\texttt{scripts/atiyah\_small\_n/}) yielded mirror-symmetric integer witnesses with $D_n < 0$ for every $26 \le n \le 34$ (\texttt{data/atiyah\_small\_n\_witnesses.json}) and stalled at $n = 26$.
The Lean certificate uses $26$ integer base points $(X_i, Y_i, Z_i)$ in $z$-monotone order:
\begin{quote}\small\raggedright
$(-65,0,0)$, $(-76,0,36)$, $(-105,0,65)$, $(-53,0,247)$, $(-23,0,994)$, $(8,0,1537)$, $(-275,-234,1778)$, $(-275,234,1778)$, $(37,-23,2189)$, $(37,23,2189)$, $(100,0,3124)$, $(26,-47,3221)$, $(26,47,3221)$, $(45,0,3314)$, $(-7,-5,3318)$, $(-7,5,3318)$, $(5,-14,3377)$, $(5,14,3377)$, $(-5,0,3390)$, $(0,-5,3412)$, $(0,5,3412)$, $(-1,0,3421)$, $(15,0,3423)$, $(0,0,3430)$, $(-385,-360,3467)$, $(-385,360,3467)$.
\end{quote}
Seven mirror pairs $(X, \pm Y, Z)$ occupy index pairs $(6,7)$, $(8,9)$, $(11,12)$, $(14,15)$, $(16,17)$, $(19,20)$, $(24,25)$; the remaining twelve points lie in $y = 0$.
Along $\mathbf{x}^{(26)}(t)_i = (\max(t,1) X_i, \max(t,1) Y_i, Z_i)$ (\texttt{x26Clamp}), $60$-digit \texttt{mpmath} gives $D_{26}(8) \approx 1.2185 \times 10^{12}$ and $D_{26}(10) \approx -3.6345 \times 10^{12}$ with root $t_0 \approx 8.289031$.
In \texttt{AtiyahSutcliffe3D26.lean}, \texttt{reflectPerm26} gives $\mathrm{Im}\, D_{26} = 0$, four affinely independent points establish \texttt{Noncoplanar3D}, and complex interval Gaussian elimination (\texttt{gaussDetCIv}) at scale $S_{26} = 10^{30}$ executed in $7$ polynomial chunks and $26$ single-pivot steps per endpoint via \texttt{decide +kernel} certifies $D_{26}(8) > 0 > D_{26}(10)$ and an IVT zero $D_{26} = 0$, proving \texttt{disproof\_atiyah\_conjecture\_1\_3d\_26} and \texttt{disproof\_atiyah\_sutcliffe\_conjecture\_2\_3d\_26}.
Combined with $n \le 4$~\cite{EastwoodNorbury2001} and $n = 5$~\cite{TanZhang2026}, the smallest $n$ for which Euclidean C1 fails lies between $6$ and $26$.

\subsection{Mirror-pair lifting: direct non-coplanar nonzero C2 disproofs at $n = 37$ and $n = 59$}
For any distinct-point planar configuration $\mathbf{x} \in C_n(\R^3)$ with $y_i \equiv 0$, define the $(n+2)$-point \emph{mirror-pair extension} $\widetilde{\mathbf{x}}_H(s) \in C_{n+2}(\R^3)$ (\texttt{mirrorPairExt}) by appending $p_n = (0, 1, H)$ and $p_{n+1} = (0, -1, H)$.
Properties~(1) and~(3) below hold generally, while property~(2) holds whenever $\mathbf{x}(s)$ contains three non-collinear points for $s \ne 0$:
\begin{enumerate}
  \item \textbf{Exact reality:} Under the transposition $\sigma = (n\;\; n+1)$, conjugation reverses $\sigma_2$ and induces row transposition on $M(\varphi)$ alongside a single sign flip on pair $(n, n+1)$ in $\mathrm{den}_{n+2}$, canceling to $\overline{D_{n+2}} = D_{n+2}$ (\texttt{atiyahDetOfSpinors\_im\_eq\_zero\_of\_lastSwap\_yRefl}).
  \item \textbf{Strict non-coplanarity:} Both $\mathbf{x}^{(35)}(s)$ and $\mathbf{x}^{(57)}(s)$ have distinct axial points on $x = 0$ and off-axis points with nonzero $x$ for $s \ne 0$; adjoining $(0, \pm 1, H)$ spans $\R^3$ (\texttt{not\_isCoplanar\_mirrorPairExt\_x35Path}).
  \item \textbf{Asymptotic decoupling:} As $H \to +\infty$, rows $n, n+1$ approach degree-$n/(n+1)$ blocks decoupling cleanly from the top-left $n \times n$ minor, giving $\lim_{H \to +\infty} D_{n+2}(\widetilde{\mathbf{x}}_H(s)) = D_n(\mathbf{x}(s))$ (\texttt{tendsto\_atiyahDet\_mirrorPairExt\_atTop}).
\end{enumerate}
Applying \texttt{mirrorPair\_anchor35} (and the $n = 57$ analogue) for a sufficiently large anchor $H > 0$ yields strictly non-coplanar 3D witnesses with $-1 < D_{37} < 0$ and $D_{37} = 0$ (\texttt{exists\_noncoplanar\_37\_violating\_conjecture\_2}, \texttt{exists\_noncoplanar\_37\_atiyahDet\_zero}), lowering our Lean-certified direct nonzero C2 counterexample in non-coplanar $\R^3$ from $n = 59$ to $n = 37$.
At concrete $H = 10^{12}$ on $\mathbf{x}^{(57)}(1)$, $60$-digit \texttt{mpmath} confirms $D_{59} \approx -0.49629962972389984453 + 6.33 \times 10^{-68} i$.

\subsection{Bounded coaxial 3D dihedral sweeps down to $n = 35$ and $\Z^3$ record $n = 46$}
Bounded random-restart Nelder--Mead sweeps using coaxial $m$-gons (\texttt{scripts/scan\_atiyah\_sutcliffe.py}, archived in \texttt{data/atiyah\_sutcliffe\_scan.json}) yield numerical observations beyond our Lean-certified theorems: a $\Z^3$ lattice witness at $n = 46$ ($K=2, m=4$, $D_{46} \approx -1.2108619818 \times 10^9 - 2.74 \times 10^{-50} i$) and sign-reversal witnesses down to $n = 35$ at $K=3, m=5$ ($D_{35} \approx -1.98163 \times 10^{16} + 5.45 \times 10^{-43} i$) and $K=2, m=7$ ($\mathrm{Re}\, D_{35} \approx -1.63886 \times 10^{12}$, $\mathrm{Im}\, D_{35} \approx 0.00940$).
Table~\ref{tab:coaxial-sweep} summarizes the smallest $n$ with $\mathrm{Re}\, D_n < 0$ across tested $(K, m)$ windows.

\begin{table}[h]
\centering
\footnotesize
\setlength{\tabcolsep}{4pt}
\caption{Smallest $n$ with $\mathrm{Re}\, D_n < 0$ and parameters $(a, b)$ found within our bounded coaxial $(K, m)$ sweeps (\texttt{data/atiyah\_sutcliffe\_scan.json}).}
\label{tab:coaxial-sweep}
\begin{tabular}{lccccc}
\toprule
Dihedral $m$ & $1$ (planar) & $2$ & $3$ & $4$ & $5$ \\
\midrule
$K = 2$ & $63\,(32,29)$ & $58\,(30,24)$ & $53\,(28,19)$ & $46\,(25,13)$ & $39\,(21,8)$ \\
$K = 3$ & $57\,(27,27)$ & -- & -- & $38\,(17,9)$ & $35\,(14,6)$ \\
\midrule
Dihedral $m$ & $6$ & $7$ & $8$ & $9$ & \\
\midrule
$K = 2$ & $36\,(18,6)$ & $35\,(15,6)$ & $37\,(16,5)$ & $38\,(15,5)$ & \\
\bottomrule
\end{tabular}
\end{table}

\section{Hyperbolic 3-Space $\Hy^3$: Disproofs Down to $n = 26$ and the $n = 3$ Five-Gap Identity}\label{sec:hyperbolic}

\subsection{Beltrami--Klein transfer at $n = 5$, $n = 35$, $n = 37$, $n = 57$, and $n = 59$}
In~\cite[\S6]{AtiyahSutcliffe2002}, Atiyah and Sutcliffe generalized their construction and conjectures to hyperbolic three-space $\Hy^3$ by replacing $v_{ij}$ with forward ideal boundary endpoints $u_{ij} \in S^2 = \partial_\infty \Hy^3$ along oriented hyperbolic geodesics.
While Eickhoff~\cite[\S7]{Eickhoff2026} established asymptotic hyperbolic singularities at large unspecified $n$, we formalize the Beltrami--Klein unit-ball model $\B^3 = \{q \in \R^3 : \|q\| < 1\}$ (\texttt{KleinBallPoints}), where geodesics are straight chords and totally geodesic hyperbolic planes equal affine planes $\cap\,\B^3$ (\texttt{IsHyperbolicCoplanar}).
For $a \ne b \in \B^3$ and $v = (b - a)/\|b - a\|$, the forward ideal endpoint is (\texttt{kleinEndpoint}):
\[
  u(a, b) = a + \tau v \in S^2, \qquad
  \tau = -(a \cdot v) + \sqrt{1 - \|a\|^2 + (a \cdot v)^2} > \|b - a\|,
\]
with $u(a, b) \ne u(b, a)$ holding universally (\texttt{kleinEndpoint\_ne\_swap}).
Independent $80$-digit evaluation confirms agreement to $< 10^{-65}$ between $D_n^{\Hy}$ in $\B^3$ and M\"obius gyroaddition in the Poincar\'e ball $\mathbb{D}^3$ ($q = 2p/(1 + \|p\|^2)$).

\begin{proposition}[Infinitesimal Euclidean recovery in $\B^3$]\label{prop:hyp-limit}
Let $\mathbf{x} \in C_n(\R^3)$. For sufficiently small $\varepsilon > 0$, we have $\varepsilon \mathbf{x} \in C_n(\B^3)$ (\texttt{scaleKlein}), $\varepsilon \mathbf{x}$ is geodesically coplanar in $\B^3$ iff $\mathbf{x}$ is affine coplanar in $\R^3$ (\texttt{isHyperbolic\allowbreak Coplanar\_scaleKlein\_iff}), and $\lim_{\varepsilon \to 0^+} D_n^{\Hy}(\varepsilon \mathbf{x}) = D_n(\mathbf{x})$ (\texttt{tendsto\_rawScaledHypDet\_zero}).
\end{proposition}

Proposition~\ref{prop:hyp-limit} yields Lean-certified hyperbolic C3 disproofs at $n = 5$ in $\Hy^2$ and $\Hy^3$ (at $\varepsilon = 10^{-7}$, $R_3^{\Hy} \approx 0.99999776100237$ in $\Hy^2$ and $0.99999776100245$ in lifted $\Hy^3$), and upgrades the Part~I hyperbolic C1/C2 bounds ($n = 57$ in $\Hy^2$, $n = 59$ in $\Hy^3$) down to \textbf{$n = 35$ in $\Hy^2$} (\texttt{exists\_planar\_35\_hyperbolic\_violating\_conjecture\_2}, \texttt{exists\_planar\_35\_hyperbolic\_atiyahDet\_zero}) and \textbf{$n = 37$ in strictly non-coplanar $\Hy^3$} (\texttt{exists\_noncoplanar\_37\_hyperbolic\_violating\_conjecture\_2}, \texttt{exists\_noncoplanar\_37\_hyperbolic\_atiyahDet\_zero}).

\subsection{General permutation reality and non-coplanar $\Hy^3$ disproof at $n = 26$}
Mirror-pair lifting in Proposition~\ref{prop:hyp-limit} required preserving reality via a single last-pair swap.
To transfer the strictly non-coplanar $n = 26$ family \texttt{x26Clamp} (which has $7$ disjoint mirror pairs) into $\Hy^3$, \texttt{AtiyahSutcliffePlanar35.lean} proves general permutation invariance \texttt{atiyahDetOfSpinors\_perm} by transposition induction alongside conjugation covariance \texttt{atiyahDetOfSpinors\_conj}, yielding the general theorem \texttt{atiyahDetOfSpinors\_im\_eq\_zero\_of\_perm\_yRefl}: any direction field $(u_{ij})$ satisfying $u_{\sigma(i),\sigma(j)} = (u_{ij,0}, -u_{ij,1}, u_{ij,2})$ for some permutation $\sigma \in S_n$ satisfies $\mathrm{Im}\, D_n(u) = 0$ identically.
Because \texttt{kleinEndpoint} commutes with $y$-reflection (\texttt{kleinEndpoint\_yRefl}), the scaled Klein family $t \mapsto \varepsilon\,\mathbf{x}^{(26)}(t) \in C_{26}(\B^3)$ satisfies $\mathrm{Im}\, D_{26}^{\Hy}(\varepsilon\,\mathbf{x}^{(26)}(t)) = 0$ on $t \in [8, 10]$ for small fixed $\varepsilon > 0$.
By Proposition~\ref{prop:hyp-limit}, $\mathrm{Re}\, D_{26}^{\Hy}(\varepsilon\,\mathbf{x}^{(26)}(8)) > 0 > \mathrm{Re}\, D_{26}^{\Hy}(\varepsilon\,\mathbf{x}^{(26)}(10))$, giving a strictly non-coplanar $26$-point hyperbolic singularity $D_{26}^{\Hy} = 0$ (\texttt{exists\_noncoplanar\_26\_hyperbolic\_atiyahDet\_zero}) and, near the root $t_0$ or endpoint $t_1$ with $-1 < D_{26}(\mathbf{x}^{(26)}(t_1)) < 0$, a direct nonzero $\Hy^3$ C2 violation $-1 < D_{26}^{\Hy} < 0$ (\texttt{exists\_noncoplanar\_26\_hyperbolic\_violating\_conjecture\_2}).

\subsection{Exact $n = 3$ $5$-gap positivity identity in $\Hy^3$}
While Hyperbolic C3 at $n = 3, 4$ remains open (for $n = 3$, C3 simplifies to $|D_3^{\Hy}| \ge 1$), \texttt{AtiyahSutcliffeNoGo.lean} establishes a general sufficient criterion $D_3^{\Hy} > 1$ from an exact algebraic difference identity.
In an equatorial hyperbolic plane $y = 0 \subset \B^3$ avoiding the stereographic south pole $(0, 0, -1)$, each ideal endpoint $u_{ij} = (u_{ij,0}, 0, u_{ij,2}) \in S^1 \subset S^2$ admits the real slope $t_{ij} := u_{ij,0} / (1 + u_{ij,2})$ and $+1$-eigenspinor $\psi_{ij} = (1, t_{ij})^\top$.
Expanding $M_3$ and $\mathrm{den}_3$ directly yields (\texttt{three\_point\_uniform\_gap\_identity}):
\begin{equation}\label{eq:five-gap}
  \det M_3(\psi) - \mathrm{den}_3(\psi) = -(t_{01} - t_{21})(t_{02} - t_{12})(t_{10} - t_{20}).
\end{equation}
Notice $\mathrm{den}_3(\psi) = (t_{10} - t_{01})(t_{20} - t_{02})(t_{21} - t_{12})$.
Whenever the six slopes obey the cyclic $5$-gap ordering
\[
  t_{20} < t_{10} < t_{12} < t_{02} < t_{01} < t_{21},
\]
all three factors of $\mathrm{den}_3(\psi)$ are strictly positive, while on the right-hand side of~\eqref{eq:five-gap} we have $t_{01} - t_{21} < 0$, $t_{02} - t_{12} > 0$, and $t_{10} - t_{20} > 0$, so the overall product is strictly positive, proving $\mathrm{Im}\, D_3^{\Hy} = 0$ and $D_3^{\Hy} > 1$ (\texttt{hyperbolic\_three\_point\_five\_gap\_gt\_one}).
Evaluating this theorem on the concrete Klein triangle $q_0 = (-3/5, 0, 0)$, $q_1 = (3/5, 0, 0)$, $q_2 = (0, 0, 4/5) \in \B^3$ (where $t_{10} = -1$, $t_{01} = 1$, and the remaining four slopes involve $\sqrt{481}/25 \in (21/25, 22/25)$) proves \texttt{exists\_hyperbolic\_three\_point\_det\_gt\_one} ($D_3^{\Hy} \approx 1.00829073076083121855 > 1$).

\section{Symplectic $\mathrm{Sp}(m)$ and Orthogonal $B_{18} = \mathrm{SO}(37)$ Disproofs}\label{sec:symplectic}

\subsection{Configuration loci and $z$-rotation parity-subblock factorization}
In Malkoun's explicit constructions for the remaining classical Lie groups~\cite{Malkoun2014,Malkoun2015}, the symplectic group $\mathrm{Sp}(m)$ ($C_m$) corresponds to centrally symmetric configurations $\mathbf{X} = (p_0, \dots, p_{2m-1}) \in C_{2m}(\R^3)$ equipped with a fixed-point-free involution $\alpha \in S_{2m}$ satisfying $p_{\alpha(i)} = -p_i$ (\texttt{SymplecticConfig m}), while the odd orthogonal group $B_m = \mathrm{SO}(2m+1)$ corresponds to $2m+1$ distinct points having a fixed origin point $p_{\mathrm{orig}} = 0$ together with $m$ antipodal pairs $p_{\alpha(i)} = -p_i$ (\texttt{OrthogonalBConfig m}).

In the $\mathrm{Sp}(m)$ setting, central symmetry pairs $(x_i, x_{\alpha(i)})$ quaternionic-conjugately, forcing $D_{2m}(\mathbf{X}) \in \R_{\ge 0}$ via Study's determinant of the quaternionic matrix~\cite{Malkoun2014}.
Consequently, unlike the $\mathrm{U}(n)$ constructions above, one cannot obtain $D_{2m} < 0$ to cross zero by a direct sign change of $D_{2m}$ itself.
Instead, \texttt{AtiyahSutcliffeSymplectic.lean} exploits the $180^\circ$ rotation around the $z$-axis, $\rho_z(x, y, z) = (-x, -y, z)$, which preserves heights $z_i$ and therefore maps the ordered stereographic slopes $w_{ij} = \frac{(x_j - x_i) + i(y_j - y_i)}{(z_j - z_i) + r_{ij}}$ (for $z_i \le z_j$) to $-w_{ij}$.
Specifically, on any $z$-monotone configuration admitting an involution $\tau \in S_N$ with $\rho_z(p_i) = p_{\tau(i)}$, the stereographic Atiyah polynomials satisfy:
\begin{itemize}
  \item $p_i(-t) = (-1)^i p_i(t)$ when $\tau(i) = i$ (so axial fixed rows vanish in columns of opposite parity modulo $2$), and
  \item $p_{\tau(i)}(t) = \mu_i\, p_i(-t)$ ($\mu_i \ne 0$) on paired rows $\tau(i) \ne i$ (multiplying coefficients by $\mu_i (-1)^k$ at degree $t^k$).
\end{itemize}
Choosing one representative row \texttt{src r} from each $\tau$-orbit reduces the full $N \times N$ determinant zero test to a single $m \times m$ real-spectrum subblock (\texttt{det\_eq\_zero\_of\_rotZ\_evenSubblock\_det\_zero} and \texttt{det\_eq\_zero\_and\_nonOrigin\_vecMul\_of\_rotZ\_oddSubblock\_det\_zero}): whenever the $m \times m$ even-degree subblock (for $\mathrm{Sp}(m)$) or odd-degree subblock (for $B_m$) has determinant zero, we lift a nonzero nullvector $v \in \C^m$ to a left nullvector $w \in \C^N$ of $M^{\mathrm{st}}$ (forcing $w_{\mathrm{orig}} = 0$ in the $B_m$ case), proving $D_N = 0$ (and, for $B_m$, linear dependence of the $2m$ non-origin rows alone).
Combining $\rho_z$ with $y$-reflection symmetry forces the $m \times m$ subblock determinant itself to be real, turning quaternionic codimension-$4$ zeros of $D_{2m}$ into robust 1D sign changes of an $m \times m$ real determinant.

\subsection{Planar $\mathrm{Sp}(20)$ ($2m = 40$) and non-coplanar 3D $\mathrm{Sp}(18)$ ($2m = 36$)}
Using outward-rounded complex interval Gaussian elimination (\texttt{gaussDetCIv}, staged via \texttt{decide +kernel} at scale $S_{\mathrm{symp}} = 10^{28}$), Lean~4 certifies two exact symplectic zero families:
\begin{itemize}
  \item \textbf{Planar $\mathrm{Sp}(20)$ ($N = 40$):} \texttt{sp20Clamp t} takes $8+8$ axial points $(0, 0, \pm(42373698 + k))$ ($0 \le k \le 7$) and $12$ antipodal planar pairs $(\pm \max(t,1) R_r, 0, \pm Z_r)$ ($R \in [1996478, 61142912]$, $Z \in [37154496, 42240438]$). Its $20 \times 20$ even-degree subblock \texttt{sp20SubMat t} flips sign on $[1, 2]$ (65/80-digit \texttt{mpmath}: $+4.8892 \times 10^{28}$ to $-3.2441 \times 10^{37}$, root $t_0 \approx 1.4319335146762459$), proving \texttt{exists\_planar\_symplectic\_20\_atiyahDet\_zero}, \texttt{disproof\_symplectic\_20\_det\_nonzero}, \texttt{disproof\_symplectic\_20\_conjecture\_1}, and \texttt{disproof\_symplectic\_20\_conjecture\_2}.
  \item \textbf{Strictly non-coplanar 3D $\mathrm{Sp}(18)$ ($N = 36$):} \texttt{sp18Clamp t} takes $6+6$ axial points $(0, 0, \pm(4408 + 5k))$ ($0 \le k \le 5$) and four $6$-point Pythagorean $(33, 56, 65)$ rings scaled horizontally by $\max(t,1)$ at $z = \pm 4043$ (inner radius $65$) and $z = \pm 4249$ (outer radius $20 \cdot 65 = 1300$). Its $18 \times 18$ even-degree subblock \texttt{sp18SubMat t} flips sign on $[7, 8]$ ($+4.8855 \times 10^{50}$ to $-3.6447 \times 10^{52}$, root $t_0 \approx 7.1955385961263573$), while the full determinant remains positive on both endpoints ($D_{36}(7) \approx 2.2855 \times 10^{24} > 0$, $D_{36}(8) \approx 4.4726 \times 10^{24} > 0$), proving \texttt{exists\_noncoplanar\_symplectic\_18\_atiyahDet\_zero} and refuting all three conjectures at $\mathrm{Sp}(18)$.
\end{itemize}
In addition, \texttt{data/atiyah\_part2\_witnesses.json} archives direct nonzero integer C2 witnesses verified in $65$-digit \texttt{mpmath} (\texttt{planar\_sp20\_40\_c2}: $D_{40} \approx +0.00158047014004968265 \in (0,1)$; \texttt{sp18\_36\_3d\_c2}: $D_{36} \approx +0.24457832515880922218 \in (0,1)$ as numerical observations complementing the certified zero theorems).

\subsection{Strictly non-coplanar 3D orthogonal $B_{18} = \mathrm{SO}(37)$ ($n = 37$)}
For odd orthogonal $B_{18}$, \texttt{b18Clamp t} places $6+6$ axial points $(0,0,\pm(5367 + 5k))$ ($0 \le k \le 5$), four $(33, 56, 65)$ Pythagorean rings at $z = \pm 5030$ (inner) and $z = \pm 5213$ (outer multiplier $19$), and the fixed origin particle $(0, 0, 0)$ at index $18$.
At scale $S_{\mathrm{symp}} = 10^{28}$, the $18 \times 18$ odd-degree subblock \texttt{b18SubMat t} flips sign from negative at $t = 8$ ($\approx -2.0568 \times 10^{48}$) to positive at $t = 9$ ($\approx +1.8592 \times 10^{51}$, root $t_0 \approx 8.7987143182361762$), giving full-determinant sign changes $D_{37}(1) \approx +4.8868 \times 10^{15}$, $D_{37}(8) \approx -6.1308 \times 10^{23}$, and $D_{37}(9) \approx +1.4426 \times 10^{25}$.
Applying \texttt{det\_eq\_zero\_and\_nonOrigin\_vecMul\_of\_rotZ\_oddSubblock\_det\_zero} yields a vanishing full determinant $D_{37} = 0$ whose left nullvector vanishes at the origin index, certifying both \texttt{exists\_noncoplanar\_orthogonalB\_18\_atiyahDet\_zero\_and\_nonOrigin\_dep} (showing that even the $36$-column non-origin polynomial family fails rank) and \texttt{disproof\_orthogonalB\_18\_det\_nonzero}, \texttt{disproof\_orthogonalB\_18\_conjecture\_1}, \texttt{disproof\_orthogonalB\_18\_conjecture\_2}.

\section{Universal $\mathrm{SU}(2) \times S_m$ Holomorphic Spinor Rigidity and Multi-Body No-Go}\label{sec:nogo}

Our counterexamples above refute Atiyah's canonical polynomial product $p_i(t)$ and its $\mathrm{Sp}(m)/B_m$ analogues.
A natural question is whether Atiyah's algebraic program can be rescued without solving nonlinear differential equations by keeping spinor-equivariant holomorphic building blocks while allowing finite sums of alternative multilinear polynomials or arbitrary $\mathrm{E}(3)$-invariant multi-body scalar weights (e.g., pairwise distances, triangles, or cross-ratios).
Section~1--2 of \texttt{AtiyahSutcliffeNoGo.lean} proves that no such modification can succeed.

\subsection{Schur--Weyl / $\mathrm{SU}(2) \times S_m$ equivariant rigidity}
Recall the homogeneous action of $U \in \mathrm{SU}(2)$ (\texttt{IsSU2}: $U U^* = I_2 \land \det U = 1$) on polynomials $P(t) = \sum_{k=0}^m a_k t^k \in \C[t]_{\le m}$ (\texttt{spinPolyAction}):
\[
  (U \cdot_m P)(t) := \sum_{k=0}^{m} a_k (U_{01} + U_{11} t)^k (U_{00} + U_{10} t)^{m-k},
\]
so that on a product $\prod_{k=0}^{m-1} (v_{k,0} + v_{k,1} t)$ this action matches transforming each spinor $v_k \mapsto U v_k$.
We call a map $\Phi : (\C^2)^m \to \C[t]$ \emph{spinor-equivariant multilinear} (\texttt{IsSpinor\allowbreak Equivariant\allowbreak Multilinear m $\Phi$}) if:
(i) $\deg \Phi(v) \le m$ everywhere,
(ii) $\Phi$ is coordinatewise $\C$-linear in each $v_k \in \C^2$,
(iii) $\Phi(v \circ \sigma) = \Phi(v)$ for every permutation $\sigma \in S_m$, and
(iv) $\Phi(U v_0, \dots, U v_{m-1}) = U \cdot_m \Phi(v)$ for every $U \in \mathrm{SU}(2)$.

\begin{proof}[Proof sketch of Theorem~\ref{thm:main-nogo}]
For the rigidity claim (\texttt{su2\_symmetric\_multilinear\_eq\_atiyahFactor}), expanding $\Phi(v)$ multilinearly in the standard basis $e_0 = (1, 0)^\top$, $e_1 = (0, 1)^\top$ of $\C^2$ and using $S_m$-symmetry reduces $\Phi$ completely to $m + 1$ polynomials $Q_r := \Phi(e_1,\dots,e_1, e_0, \dots, e_0)$ with $r$ leading copies of $e_1$ ($0 \le r \le m$).
Applying the diagonal matrix $U_\zeta = \mathrm{diag}(\zeta, \bar{\zeta}) \in \mathrm{SU}(2)$ with Pythagorean unit $\zeta = (3 + 4i)/5$ forces $(Q_r)_{t^s} = 0$ for all $0 \le r, s \le m$ with $s \ne r$, since $\zeta^2 = (-7 + 24i)/25$ is not a root of unity so $\bar{\zeta}^s \zeta^{m-s} \ne \bar{\zeta}^r \zeta^{m-r}$; hence $Q_r(t) = c_r t^r$.
Applying the Hadamard unitary
\[
  U_H = \frac{1}{\sqrt{2}}\begin{pmatrix} 1 & -1 \\ 1 & 1 \end{pmatrix} \in \mathrm{SU}(2)
\]
to $(e_0, \dots, e_0)$ equates $\sum_{r=0}^m \binom{m}{r} 2^{-m/2} c_r t^r$ with $2^{-m/2} c_0 (t + 1)^m = \sum_{r=0}^m \binom{m}{r} 2^{-m/2} c_0 t^r$, forcing $c_r = c_0$ for all $0 \le r \le m$.
Matching subset expansions establishes $\Phi(v) = c_0 \prod_{k=0}^{m-1} (v_{k,0} + v_{k,1} t)$ identically.
\end{proof}

\subsection{Collapse of multi-body scalar-weighted ensembles}
For any finite index set $\iota$, any family $(\Phi_\alpha)_{\alpha \in \iota}$ satisfying \texttt{IsSpinor\allowbreak Equivariant\allowbreak Multilinear (n-1) ($\Phi_\alpha$)}, and completely arbitrary complex scalar weights $W_\alpha : C_n(\R^3) \times \{0, \dots, n - 1\} \to \C$, substituting the canonical pairwise spinors $v_{i,r} = \psi^{\mathrm{can}}_{i, \mathrm{omit}_i(r)}$ into $P_i(\mathbf{x}; t) = \sum_{\alpha \in \iota} W_\alpha(\mathbf{x}, i) \Phi_\alpha((v_{i,r})_r)$ (\texttt{weightedSpinorCandidate}) yields
\[
  P_i(\mathbf{x}; t) = \mu_i(\mathbf{x})\, p_i^{\mathrm{Atiyah}}(t), \qquad \mu_i(\mathbf{x}) := \sum_{\alpha \in \iota} W_\alpha(\mathbf{x}, i) c_\alpha,
\]
independently of $t$ (\texttt{weightedSpinorCandidate\_eq\_smul\_atiyahPoly}).
Scaling the rows of a linearly dependent family of polynomials by arbitrary scalars $\mu_i \in \C$ (whether zero or nonzero) preserves linear dependence (\texttt{not\_linearIndependent\_scale\_of\_not\_linearIndependent}).
Evaluating at our certified singular configurations produces the three unconditional No-Go theorems at $n = 26$ ($\mathbb{R}^3$), $n = 35$ ($\mathbb{R}^2$), and $2m = 36$ ($\mathrm{Sp}(18)$).

\begin{remark}[Holomorphic multilinearity vs.\ Nahm's equations]
While smooth $\mathrm{SO}(3) \times S_n$-equivariant Berry--Robbins maps exist unconditionally via solutions to Nahm's equations~\cite{AtiyahBielawski2002}, Theorem~\ref{thm:main-nogo} explains why no elementary holomorphic spinor-polynomial ansatz---even with arbitrary $\mathrm{E}(3)$-invariant multi-body distance weights---can escape the $n = 26, 35, 36$ obstructions without breaking either $\mathrm{SU}(2)$-equivariant holomorphicity or coordinatewise multilinearity in the pairwise directions.
\end{remark}

\section{Formal Verification in Lean~4}

All formalized proofs reside in the seven Track~B Lean~4 modules (together with shared preliminaries in \texttt{MalkounCounterexample.lean}):
\begin{quote}\small\raggedright
\texttt{AtiyahSutcliffe.lean}, \texttt{AtiyahSutcliffeExtensions.lean}, \texttt{AtiyahSutcliffe3D36.lean}, \texttt{AtiyahSutcliffe3D26.lean}, \texttt{AtiyahSutcliffePlanar35.lean}, \texttt{AtiyahSutcliffeSymplectic.lean}, \texttt{AtiyahSutcliffeNoGo.lean}.
\end{quote}
All modules are imported by \texttt{lean/Matematyka.lean} and compiled against Lean~\texttt{v4.37.0-rc1} and Mathlib (\texttt{lake build}).
A standalone reproduction repository containing all Lean~4 proofs, Python/\texttt{mpmath} scripts, witness datasets, and this manuscript is archived at \url{https://github.com/tomasztarlowski/atiyah-sutcliffe-counterexamples} and permanently preserved on Zenodo at \url{https://doi.org/10.5281/zenodo.24443573}.
All modules contain zero occurrences of \texttt{sorry}, \texttt{admit}, \texttt{native\_decide}, or custom axioms (staged table identities use kernel-only \texttt{decide +kernel}), compile cleanly with zero warnings, and automated \texttt{\#print axioms} auditing confirms that every exported theorem depends exclusively on \texttt{propext}, \texttt{Classical.choice}, and \texttt{Quot.sound}.
Tables~\ref{tab:lean-theorems-1} and~\ref{tab:lean-theorems-2} list the principal declarations verified across Parts~I and~II.

\begin{table}[htbp]
\centering
\scriptsize
\setlength{\tabcolsep}{4pt}
\renewcommand{\arraystretch}{0.92}
\caption{Principal Lean~4 declarations verified in \texttt{AtiyahSutcliffe}, \texttt{AtiyahSutcliffeExtensions}, \texttt{AtiyahSutcliffe3D36}, and \texttt{AtiyahSutcliffe3D26}.}
\label{tab:lean-theorems-1}
\begin{tabular}{>{\ttfamily\raggedright\arraybackslash}p{0.46\textwidth}>{\raggedright\arraybackslash}p{0.47\textwidth}}
\toprule
\textrm{Lean Declaration} & Statement Formalized \\
\midrule
exists\_five\_planar\_points\_violating\_conjecture\_3 & Existential $x \in C_5(\R^3)$ violating C3 (witnessed by planar integer $\mathbf{x}^{(5)}$) \\
disproof\_atiyah\_sutcliffe\_conjecture\_3 & Negation $\neg \mathrm{Atiyah\allowbreak Sutcliffe\allowbreak Conjecture3Prop}$ ($n \ge 3$) \\
exists\_five\_noncoplanar\_points\_violating\_conjecture\_3 & Existential non-coplanar ($\neg \mathrm{IsCoplanar}$) $n = 5$ C3 witness \\
x57Path\_one\_re\_neg\_and\_norm\_lt\_one & Explicit integer $\mathbf{x}^{(57)}(1)$ satisfies $-1 < D_{57} < 0$ ($S_{57}=10^{40}$) \\
exists\_planar\_57\_atiyahDet\_zero & Existential $x \in C_{57}(\R^3)$ with $D_{57} = 0$ (IVT along $\mathbf{x}^{(57)}$) \\
x63Path\_one\_norm\_lt\_one & Explicit integer $\mathbf{x}^{(63)}(1)$ satisfies $|D_{63}| < 1$ ($S_{63}=10^{35}$) \\
det\_stereoMat\_eq\_even\_odd & Antipodal parity factorization $36 \to 18 \times 18$ (\texttt{x36Clamp}) \\
exists\_noncoplanar\_3d\_36\_atiyahDet\_zero & Non-coplanar \texttt{Noncoplanar3D} $n = 36$ with $D_{36} = 0$ \\
disproof\_atiyah\_conjecture\_1\_3d\_36 & Non-coplanar 3D $n = 36$ refutation $\neg \mathrm{Atiyah\allowbreak Conjecture1Prop}$ \\
gaussDetCIv\_mem & Soundness of certified complex interval Gaussian elimination \\
exists\_noncoplanar\_3d\_26\_atiyahDet\_neg & Non-coplanar \texttt{Noncoplanar3D} $n = 26$ with $\mathrm{Re}\, D_{26} < 0$, $\mathrm{Im}\, D_{26} = 0$ \\
exists\_noncoplanar\_3d\_26\_atiyahDet\_zero & Non-coplanar \texttt{Noncoplanar3D} $n = 26$ with $D_{26} = 0$ (IVT on $[8, 10]$) \\
disproof\_atiyah\_conjecture\_1\_3d\_26 & Non-coplanar 3D $n = 26$ refutation $\neg \mathrm{Atiyah\allowbreak Conjecture1Prop}$ \\
disproof\_atiyah\_sutcliffe\_conjecture\_2\_3d\_26 & Non-coplanar 3D $n = 26$ refutation $\neg \mathrm{Atiyah\allowbreak Sutcliffe\allowbreak Conjecture2Prop}$ \\
exists\_five\_planar\_hyperbolic\_violating\_conjecture\_3 & Existential $\Hy^2$ $n = 5$ witness violating Hyperbolic C3 \\
exists\_five\_noncoplanar\_hyperbolic\_violating\_conjecture\_3 & Non-coplanar $\Hy^3$ $n = 5$ witness violating Hyperbolic C3 \\
disproof\_hyperbolic\_atiyah\_sutcliffe\_conjecture\_3 & Negation $\neg \mathrm{Hyperbolic\allowbreak Atiyah\allowbreak Sutcliffe\allowbreak Conjecture3Prop}$ \\
exists\_noncoplanar\_59\_violating\_conjecture\_2 & Non-coplanar $n = 59$ witness with $-1 < D_{59} < 0$ \\
disproof\_hyperbolic\_atiyah\_conjecture\_1 & Negation $\neg \mathrm{Hyperbolic\allowbreak Atiyah\allowbreak Conjecture1Prop}$ \\
disproof\_hyperbolic\_atiyah\_sutcliffe\_conjecture\_2 & Negation $\neg \mathrm{Hyperbolic\allowbreak Atiyah\allowbreak Sutcliffe\allowbreak Conjecture2Prop}$ \\
\bottomrule
\end{tabular}
\end{table}

\begin{table}[htbp]
\centering
\scriptsize
\setlength{\tabcolsep}{4pt}
\renewcommand{\arraystretch}{0.92}
\caption{Principal Lean~4 declarations verified in \texttt{AtiyahSutcliffePlanar35}, \texttt{AtiyahSutcliffeSymplectic}, and \texttt{AtiyahSutcliffeNoGo}.}
\label{tab:lean-theorems-2}
\begin{tabular}{>{\ttfamily\raggedright\arraybackslash}p{0.49\textwidth}>{\raggedright\arraybackslash}p{0.44\textwidth}}
\toprule
\textrm{Lean Declaration} & Statement Formalized \\
\midrule
x35Path\_one\_re\_neg\_and\_norm\_lt\_one & Direct integer endpoint $-1 < D_{35}(\mathbf{x}^{(35)}(1)) < 0$ ($S_{35}=10^{34}$) \\
exists\_planar\_35\_violating\_conjecture\_2 & Planar $n = 35$ witness violating C2 ($-1 < D_{35} < 0$) \\
exists\_planar\_35\_atiyahDet\_zero & Planar $n = 35$ singular witness ($D_{35} = 0$, IVT along \texttt{x35Path}) \\
exists\_noncoplanar\_37\_violating\_conjecture\_2 & Non-coplanar 3D $n = 37$ mirror-pair witness ($-1 < D_{37} < 0$) \\
exists\_planar\_35\_hyperbolic\_violating\_conjecture\_2 & Planar $\Hy^2$ $n = 35$ witness with $-1 < D_{35}^{\Hy} < 0$ \\
exists\_noncoplanar\_37\_hyperbolic\_violating\_conjecture\_2 & Non-coplanar $\Hy^3$ $n = 37$ witness with $-1 < D_{37}^{\Hy} < 0$ \\
atiyahDetOfSpinors\_im\_eq\_zero\_of\_perm\_yRefl & Permutation $y$-reflection reality theorem for general spinor tables \\
exists\_noncoplanar\_26\_hyperbolic\_atiyahDet\_zero & Non-coplanar $\Hy^3$ $n = 26$ singular witness ($D_{26}^{\Hy} = 0$) \\
exists\_noncoplanar\_26\_hyperbolic\_violating\_conjecture\_2 & Non-coplanar $\Hy^3$ $n = 26$ direct C2 witness ($-1 < D_{26}^{\Hy} < 0$) \\
exists\_planar\_symplectic\_20\_atiyahDet\_zero & Planar $\mathrm{Sp}(20)$ ($N = 40$) singular configuration ($D_{40} = 0$) \\
disproof\_symplectic\_20\_conjecture\_1 & Refutation $\neg \mathrm{Symplectic\allowbreak Atiyah\allowbreak Conjecture1Prop}\; 20$ (plus \texttt{det\_nonzero}, \texttt{C2}) \\
exists\_noncoplanar\_symplectic\_18\_atiyahDet\_zero & Non-coplanar 3D $\mathrm{Sp}(18)$ ($N = 36$) singular witness ($D_{36} = 0$) \\
disproof\_symplectic\_18\_conjecture\_1 & Refutation $\neg \mathrm{Symplectic\allowbreak Atiyah\allowbreak Conjecture1Prop}\; 18$ (plus \texttt{det\_nonzero}, \texttt{C2}) \\
exists\_noncoplanar\_orthogonalB\_18\_atiyahDet\_zero\_and\_nonOrigin\_dep & Non-coplanar $B_{18}$ ($N = 37$) $D_{37} = 0$ with $36$ non-origin rows dependent \\
disproof\_orthogonalB\_18\_conjecture\_1 & Refutation $\neg \mathrm{OrthogonalB\allowbreak Atiyah\allowbreak Conjecture1Prop}\; 18$ (plus \texttt{det\_nonzero}, \texttt{C2}) \\
su2\_symmetric\_multilinear\_eq\_atiyahFactor & $\mathrm{SU}(2) \times S_m$ multilinear spinor polynomial rigidity \\
weightedSpinorCandidate\_eq\_smul\_atiyahPoly & Row-wise collapse of arbitrary scalar-weighted ensembles \\
universal\_holomorphic\_spinor\_nogo\_noncoplanar\_3d\_26 & Universal weighted spinor No-Go at non-coplanar $n = 26$ \\
universal\_holomorphic\_spinor\_nogo\_planar\_35 & Universal weighted spinor No-Go at planar $n = 35$ \\
universal\_holomorphic\_spinor\_nogo\_symplectic\_36 & Universal weighted spinor No-Go at non-coplanar $\mathrm{Sp}(18)$ ($n = 36$) \\
three\_point\_uniform\_gap\_identity & Exact algebraic identity $\det M_3 - \mathrm{den}_3$ in uniform chart \\
hyperbolic\_three\_point\_five\_gap\_gt\_one & Cyclic $5$-gap positivity theorem $D_3^{\Hy} > 1$ in $\B^3$ \\
exists\_hyperbolic\_three\_point\_det\_gt\_one & Concrete Klein triangle witness with $D_3^{\Hy} > 1$ \\
\bottomrule
\end{tabular}
\end{table}

\clearpage
\subsection*{Acknowledgments and AI Disclosure}
Computational witness discovery, Lean~4 proof engineering, and manuscript drafting were carried out in an interactive human--AI research workflow using Gemini, Claude, and OpenAI Codex models.
All formal theorems are machine-checked by the Lean~4 kernel (\texttt{v4.37.0-rc1} + Mathlib) with zero \texttt{sorry}, zero \texttt{native\_decide}, and standard axioms \texttt{\{propext, Classical.choice, Quot.sound\}} only, and all reported numerical quantities were independently verified in $60$-, $65$-, and $80$-digit \texttt{mpmath} arithmetic.

\begin{thebibliography}{15}

\bibitem{Akhtyamov2026}
A.~Akhtyamov,
\emph{Formal disproof of Atiyah--Sutcliffe Conjecture~1 in Lean~4},
GitHub repository \url{https://github.com/azakhtyamov/atiyah-sutcliffe-disproof} and pull request \texttt{google-deepmind/formal-conjectures\#6966} (October 2026).

\bibitem{Atiyah2001}
M.~F.~Atiyah,
\emph{Configurations of points},
Phil. Trans. R. Soc. Lond. A \textbf{359} (2001), 1375--1387.

\bibitem{AtiyahBielawski2002}
M.~F.~Atiyah and R.~Bielawski,
\emph{Nahm's equations, configuration spaces and flag manifolds},
Bull. Braz. Math. Soc. (N.S.) \textbf{33} (2002), 157--176 (\texttt{arXiv:math/0110112}).

\bibitem{AtiyahSutcliffe2002}
M.~F.~Atiyah and P.~Sutcliffe,
\emph{The geometry of point particles},
Proc. R. Soc. Lond. A \textbf{458} (2002), 1089--1115.

\bibitem{BerryRobbins1997}
M.~V.~Berry and J.~M.~Robbins,
\emph{Indistinguishability for quantum particles: spin, statistics and the geometric phase},
Proc. R. Soc. Lond. A \textbf{453} (1997), 1771--1790.

\bibitem{BouKhuzamJohnson2014}
M.~N.~Bou~Khuzam and M.~J.~Johnson,
\emph{On the conjectures regarding the $4$-point Atiyah determinant},
SIGMA \textbf{10} (2014), 070, 9 pages (\texttt{arXiv:1401.2787}).

\bibitem{Djokovic2002}
D.~\v{Z}.~\DJ okovi\'c,
\emph{Proof of Atiyah's conjecture for two special types of configurations},
Electron. J. Linear Algebra \textbf{9} (2002), 132--137 (\texttt{arXiv:math/0205221}).

\bibitem{EastwoodNorbury2001}
M.~G.~Eastwood and P.~Norbury,
\emph{A proof of Atiyah's conjecture on configurations of four points in Euclidean three-space},
Geom. Topol. \textbf{5} (2001), 885--893 (\texttt{arXiv:math/0109161}).

\bibitem{Eickhoff2026}
J.~Eickhoff,
\emph{Singular Atiyah Configurations},
preprint \texttt{arXiv:2610.33090} (2026).

\bibitem{Malkoun2014}
J.~Malkoun,
\emph{Configurations of points and the symplectic Berry--Robbins problem},
SIGMA \textbf{10} (2014), 112, 6 pages (\texttt{arXiv:1407.8291}).

\bibitem{Malkoun2015}
J.~Malkoun,
\emph{Constructions of $\mathrm{SU}(2)$ and Weyl equivariant maps for all classical groups},
preprint \texttt{arXiv:1508.04076} (2015).

\bibitem{Malkoun2025}
J.~Malkoun,
\emph{Finite graphs and configurations of points},
preprint \texttt{arXiv:2508.13472} (2025).

\bibitem{MazurPetrenko2011}
M.~Mazur and B.~V.~Petrenko,
\emph{On the conjectures of Atiyah and Sutcliffe},
Geom. Dedicata \textbf{158} (2012), 329--342 (\texttt{arXiv:1102.4662}).

\bibitem{TanZhang2026}
H.~Tan and Y.~Zhang,
\emph{The Atiyah--Sutcliffe determinant at five points: positivity and failure of the deletion inequality},
preprint, Zenodo (4 October 2026), \url{https://doi.org/10.5281/zenodo.23142909}.

\end{thebibliography}

\end{document}
